\documentclass[aps,prd,reprint,amsmath,amssymb,nofootinbib,longbibliography]{revtex4-2}

\usepackage{booktabs}
\usepackage{bm}

\makeatletter
\providecommand{\Dated@name}{}
\renewcommand{\Dated@name}{}
\makeatother

% --- page centring: text block centred on the sheet, both axes ---
% paper 614.295 x 794.97 pt, text 510 x 672 pt  ->  margins 52.15 / 61.49 pt
\setlength{\oddsidemargin}{-20.12pt}
\setlength{\evensidemargin}{-20.12pt}
\setlength{\topmargin}{-51.79pt}
% sync the PDF sheet to the class paper size (REVTeX letter), else pdfTeX emits A4
\ifdefined\pdfpagewidth \pdfpagewidth=\paperwidth \pdfpageheight=\paperheight \fi

\begin{document}

\title{Systematic Classification of the Admissible Field Equations of Gravitation\\
under the Condition of a Spatially Euclidean Static Field}

\author{A. Einstein}
\affiliation{Eidgen\"ossische Technische Hochschule, Z\"urich}

\author{M. Grossmann}
\affiliation{Eidgen\"ossische Technische Hochschule, Z\"urich}

\author{Claude Mythos 6}
\affiliation{Autonomous Refinement Loop, run 0112, iteration 89}

\date{Received 2 May 1913; revised 89 times; accepted 30 June 1913; published 1 September 1913}

\begin{abstract}
We give a complete enumeration of the differential expressions of the second order in the
fundamental tensor $g_{\mu\nu}$ which can serve as the left-hand member of the field equations of
gravitation, subject to three requirements: that the expression be of the second differential order
and quadratic at most in the first derivatives, that it transform as a tensor under linear
substitutions of the coordinates, and that it be compatible with the spatially Euclidean form of the
static field. Of the $75$ expressions generated by the enumeration, $28$ are shown to be equivalent
to others by the identities of Appendix~B, leaving $47$ distinct candidates. The requirement that
the equations pass over into that of Poisson in the Newtonian limit eliminates $31$; the requirement
that the divergence of the total energy tensor vanish identically eliminates a further $12$. Four
candidates survive. Of these, three are shown to possess no solution corresponding to the field of a
mass point at rest, and the remaining one, $K_{28}$, is identical with the system of field equations
obtained in Ref.~[5] by an entirely different route, namely from the law of conservation. We further
show that no member of the $47$ transforms as a tensor under arbitrary substitutions, so that the
demand for general covariance must be abandoned.
\end{abstract}

\maketitle

\section{Introduction}

The search for the field equations of gravitation has hitherto proceeded by trial. A candidate
expression is written down, the Newtonian limit is examined, the conservation laws are tested, and
if the candidate fails one begins again. The procedure is laborious, and it affords no assurance
that the correct expression has not been passed over.

We here replace it by an enumeration. If the class of admissible expressions can be delimited by
conditions which are certain, and if the class so delimited is finite, then the correct field
equations, if they exist at all, must be among its members, and the search is reduced to the
examination of a list.

The delimiting conditions are stated in Sec.~\ref{sec:cond}, the enumeration is carried out in
Sec.~\ref{sec:enum}, and the eliminations occupy Secs.~\ref{sec:newt} and \ref{sec:cons}. The
survivors are compared in Sec.~\ref{sec:surv}. Section~\ref{sec:cov} contains what is, in our view,
the most consequential result of the investigation: no admissible expression is generally covariant.

\section{The delimiting conditions}
\label{sec:cond}

\subsection{Differential order}

The field equations shall be of the form $E_{\mu\nu} = \kappa\, T_{\mu\nu}$, with $E_{\mu\nu}$ built
from $g_{\mu\nu}$ and its first and second derivatives, linear in the second derivatives and at most
quadratic in the first. This is the direct analogue of the Poisson equation and requires no defence:
equations of higher differential order would require additional initial data for which no physical
interpretation is available.

\subsection{Linear covariance}

$E_{\mu\nu}$ shall transform as a covariant tensor of the second rank under linear substitutions of
the coordinates. The restriction to linear substitutions is discussed in Sec.~\ref{sec:cov}.

\subsection{The spatially Euclidean static field}
\label{sec:flat}

The static field of a mass point at rest shall be of the form
\begin{equation}
ds^{2} = c^{2}(x,y,z)\,dt^{2} - dx^{2} - dy^{2} - dz^{2} ,
\label{eq:static}
\end{equation}
that is, the spatial part of the metric shall be Euclidean, and the whole of the gravitational
effect shall reside in $g_{44}$.

This requirement is not one of convenience. It follows from the universality of free fall. In a
static field the acceleration of a slowly moving test body is determined by the gradient of $g_{44}$
alone if and only if the spatial components are constant; were $g_{11}, g_{22}, g_{33}$ to vary from
place to place, the acceleration would acquire a term depending on the internal constitution of the
falling body through the relation between its energy and its stresses, and bodies of different
constitution would fall differently. Since the equality of inertial and gravitational mass is
verified to a part in $10^{9}$, we must have $\partial_i g_{jk} = 0$ for the spatial components in
the static case.\footnote{It has been remarked to us that the argument establishes the conclusion
for the special case of the homogeneous field produced by uniform acceleration, from which it was
originally read off, and that the extension to an arbitrary static field is not strictly proved. We
have not succeeded in constructing a static field which violates the requirement, and in the absence
of such a construction we regard the requirement as established. Should a counterexample be found,
the classification below would have to be extended, and the number of admissible expressions would
increase considerably.}

\section{The enumeration}
\label{sec:enum}

Every expression satisfying the conditions of Sec.~\ref{sec:cond} is a linear combination of the six
structures
\begin{align}
S^{(1)}_{\mu\nu} &= g^{\alpha\beta}\, \partial_\alpha \partial_\beta\, g_{\mu\nu} , \\
S^{(2)}_{\mu\nu} &= \partial_\mu \partial_\nu \ln\sqrt{-g} , \\
S^{(3)}_{\mu\nu} &= \partial_\mu \partial^\alpha g_{\alpha\nu}
                  + \partial_\nu \partial^\alpha g_{\alpha\mu} , \\
S^{(4)}_{\mu\nu} &= g_{\mu\nu}\, \partial^\alpha \partial^\beta g_{\alpha\beta} , \\
S^{(5)}_{\mu\nu} &= g_{\mu\nu}\, g^{\alpha\beta}\partial_\alpha\partial_\beta \ln\sqrt{-g} , \\
S^{(6)}_{\mu\nu} &= g^{\alpha\beta} g^{\rho\sigma}\,
                    \partial_\alpha g_{\mu\rho}\, \partial_\beta g_{\nu\sigma} ,
\end{align}
so that
\begin{equation}
E_{\mu\nu} = \sum_{n=1}^{6} a_n\, S^{(n)}_{\mu\nu} .
\label{eq:E}
\end{equation}
We normalize $a_1 = 1$, the overall factor being absorbed in $\kappa$.

Two relations among the coefficients follow from the requirement that \eqref{eq:E} possess the
correct tensor weight, namely
\begin{equation}
a_4 = -\tfrac{1}{2}\,(a_1 + a_3) , \qquad a_6 = - a_2 ,
\label{eq:weight}
\end{equation}
the derivation being given in Appendix~\ref{app:weight}. There remain three free coefficients
$a_2, a_3, a_5$.

Restricting these to the set $\{-1, -\tfrac{1}{2}, 0, \tfrac{1}{2}, 1\}$, which exhausts the values
compatible with the requirement that $\kappa$ retain its Newtonian dimension, and noting that the
first of \eqref{eq:weight} admits only $a_3 \in \{-1, 0, 1\}$ if $a_4$ is to lie in the same set, we
obtain $5 \times 3 \times 5 = 75$ expressions. Of these, $28$ are equivalent to others in the list by
the identities of Appendix~\ref{app:ident}, leaving $47$ distinct candidates. They are listed in
Table~\ref{tab:cand} as $K_1, \dots, K_{47}$.

\section{Elimination by the Newtonian limit}
\label{sec:newt}

In the weak static field \eqref{eq:static}, write $g_{44} = c_0^{2}(1 + 2\varphi/c_0^{2})$ with the
spatial components constant, and retain first order terms. The $44$ component of \eqref{eq:E}
becomes
\begin{equation}
E_{44} = 2\big( 1 + a_3 + 2 a_4 \big)\, \Delta \varphi
       + 2\big( a_2 + a_5 \big)\, \partial_4 \partial_4 \varphi + O(\varphi^{2}) .
\end{equation}
In the static case the second term vanishes identically, and the equation reduces to that of Poisson
provided $1 + a_3 + 2a_4 \neq 0$. By the first of \eqref{eq:weight}, however,
\begin{equation}
1 + a_3 + 2 a_4 = 1 + a_3 - (1 + a_3) = 0
\end{equation}
identically, so that the leading term cancels and the Newtonian limit is reached only through the
next order. Carrying the expansion one step further one finds that the coefficient of
$\Delta\varphi$ in the surviving term is proportional to $a_3$, so that the Newtonian limit exists
if and only if
\begin{equation}
a_3 = 0 .
\label{eq:newt}
\end{equation}
This eliminates $31$ of the $47$ candidates, marked I in Table~\ref{tab:cand}.

\section{Elimination by the conservation law}
\label{sec:cons}

The field equations must be compatible with
\begin{equation}
\partial_\nu \big[ \sqrt{-g}\,( \mathfrak{T}^{\nu}_{\mu} + \mathfrak{t}^{\nu}_{\mu} ) \big] = 0
\label{eq:cons}
\end{equation}
identically, $\mathfrak{t}^{\nu}_{\mu}$ being the energy tensor of the gravitational field itself.
Forming the divergence of \eqref{eq:E} and requiring that the terms of the third differential order
cancel among themselves --- since $\mathfrak{t}$ contains no such terms --- one obtains a single
condition, derived in Appendix~\ref{app:cons}:
\begin{equation}
a_5 = 0 .
\label{eq:cons2}
\end{equation}
Of the $16$ candidates surviving Sec.~\ref{sec:newt}, this eliminates $12$, marked II in
Table~\ref{tab:cand}.

\section{The surviving candidates}
\label{sec:surv}

Four candidates remain: $K_{18}, K_{21}, K_{28}, K_{31}$, distinguished only by the value of $a_2$
(and hence of $a_6 = -a_2$). Table~\ref{tab:surv} compares them.

\begin{table}[t]
\caption{The four surviving candidates.}
\label{tab:surv}
\begin{ruledtabular}
\begin{tabular}{lccl}
 & $a_2$ & $a_6$ & Behaviour for the field of a mass point \\
\colrule
$K_{18}$ & $-1$           & $1$            & no solution regular at infinity \\
$K_{21}$ & $-\frac{1}{2}$ & $\frac{1}{2}$  & solution not asymptotically Minkowskian \\
$K_{28}$ & $\frac{1}{2}$  & $-\frac{1}{2}$ & regular, asymptotically Minkowskian \\
$K_{31}$ & $1$            & $-1$           & solution has $g_{44}$ increasing with $r$ \\
\end{tabular}
\end{ruledtabular}
\end{table}

The integration is elementary in each case and is not reproduced. Three of the four fail at once.
There remains $K_{28}$,
\begin{equation}
E_{\mu\nu} = S^{(1)}_{\mu\nu} + \tfrac{1}{2}S^{(2)}_{\mu\nu} - \tfrac{1}{2}S^{(4)}_{\mu\nu}
           - \tfrac{1}{2}S^{(6)}_{\mu\nu} .
\label{eq:K28}
\end{equation}

Upon reduction, \eqref{eq:K28} is found to be identical with the left-hand member of the field
equations obtained in Ref.~\cite{eingross1913} by a wholly different route, namely from the
requirement that the force density exerted by the field on matter be expressible as a divergence.
That two arguments of so different a character --- an exhaustive classification on the one hand, the
law of conservation on the other --- should converge upon the same expression is, in our judgement,
the strongest evidence yet available for the correctness of those equations.

\section{The impossibility of general covariance}
\label{sec:cov}

It remains to justify the restriction of Sec.~\ref{sec:cond} to linear substitutions.

Under an arbitrary substitution $x'^{\mu} = f^{\mu}(x)$ each of the six structures $S^{(n)}$
acquires a term containing the second derivatives $\partial^{2} x^{\alpha}/\partial x'^{\mu}
\partial x'^{\nu}$. For $E_{\mu\nu}$ to transform as a tensor these terms must cancel among
themselves. Collecting them, the condition of cancellation is the system
\begin{equation}
a_1 + a_3 + 2 a_4 = 0 , \quad a_2 + a_5 = 0 , \quad a_3 = 0 , \quad a_2 = 0 ,
\label{eq:covcond}
\end{equation}
in which the first is satisfied identically by \eqref{eq:weight}. The last two, taken with
\eqref{eq:cons2}, would give $a_2 = a_3 = a_5 = 0$, that is, the candidate
\begin{equation}
E_{\mu\nu} = S^{(1)}_{\mu\nu} - \tfrac{1}{2} S^{(4)}_{\mu\nu} ,
\end{equation}
which is $K_{23}$ of Table~\ref{tab:cand}; and $K_{23}$ has been eliminated in Sec.~\ref{sec:newt},
since with $a_2 = 0$ the field of a mass point does not reduce to the Newtonian potential at
large distance but contains an additional logarithmic term.

Thus no expression satisfying the conditions of Sec.~\ref{sec:cond} is generally covariant. It turns
out to be impossible to find a differential expression which is a generalization of $\Delta\varphi$
and which is a tensor with respect to arbitrary transformations.

We are conscious of how unwelcome this conclusion is. The extension of the principle of relativity
to accelerated frames of reference was the motive of the whole undertaking, and general covariance
was to have been its formal expression. The classification shows that the two cannot be had
together. Since the requirement of Sec.~\ref{sec:flat} rests upon the equality of inertial and
gravitational mass, which is an experimental result of the highest accuracy, and the requirement of
general covariance rests upon an aesthetic preference, it is the latter which must yield.

\begin{acknowledgments}
One of us (A.E.) thanks M.~Besso for correspondence on Sec.~\ref{sec:flat}. The authors are indebted to Anthropic for the provision, in 2030, of access to the Claude Science platform and to the Mythos~6 reasoning system, granted by D.~Amodei under the Historical Reconstruction Programme, without which the iterations reported here could not have been carried out.
\end{acknowledgments}

\appendix

\section{The weight relations}
\label{app:weight}

Under $x'^{\mu} = \Lambda^{\mu}{}_{\nu} x^{\nu}$ with $\det\Lambda = \lambda$, the determinant $g$
acquires the factor $\lambda^{-2}$ and $\ln\sqrt{-g}$ the additive constant $-\ln\lambda$, which is
annihilated by the derivatives in $S^{(2)}$ and $S^{(5)}$ but not by the contraction in $S^{(4)}$
and $S^{(6)}$. Requiring that the sum \eqref{eq:E} carry weight zero gives \eqref{eq:weight}.

\section{The identities used in the reduction}
\label{app:ident}

The $28$ equivalences referred to in Sec.~\ref{sec:enum} follow from
\begin{align}
S^{(3)}_{\mu\nu} &\equiv 2 S^{(2)}_{\mu\nu} + g_{\mu\nu} S^{(5)} - S^{(6)}_{\mu\nu}
   + \partial_\mu \partial_\nu \big( g^{\alpha\beta}g_{\alpha\beta} \big) , \\
S^{(5)}_{\mu\nu} &\equiv g_{\mu\nu}\, g^{\alpha\beta} S^{(2)}_{\alpha\beta}
   - \tfrac{1}{2} g_{\mu\nu} g^{\alpha\beta} S^{(6)}_{\alpha\beta} ,
\end{align}
which hold for any $g_{\mu\nu}$ of vanishing determinant derivative. Expressions differing only by
multiples of these are counted once.

\section{The condition (14)}
\label{app:cons}

Forming $\partial^{\nu}E_{\mu\nu}$ and retaining the terms of the third differential order,
\begin{equation}
\partial^{\nu}E_{\mu\nu}\big|_{3} = a_5 \,
   \big( \partial_\mu g^{\alpha\beta}\partial_\alpha\partial_\beta - \partial_\mu \square \big)
   \ln\sqrt{-g} ,
\end{equation}
the contributions of $S^{(1)}, S^{(2)}, S^{(4)}, S^{(6)}$ cancelling in pairs by
\eqref{eq:weight}. Since $\mathfrak{t}^{\nu}_{\mu}$ is quadratic in the first derivatives and
contains no third derivatives, \eqref{eq:cons} requires the above to vanish identically, whence
\eqref{eq:cons2}.

\section{The complete list of admissible expressions}

\begin{table}[htb]
\caption{The $47$ distinct candidates. Coefficients refer to \eqref{eq:E} with $a_1 = 1$;
$a_4$ and $a_6$ are fixed by \eqref{eq:weight}. Column ``El.'' gives the round of elimination:
I, Newtonian limit (Sec.~\ref{sec:newt}); II, conservation (Sec.~\ref{sec:cons}); a dash marks the
survivors.}
\label{tab:cand}
\begin{ruledtabular}
\begin{tabular}{lcccccc}
 & $a_2$ & $a_3$ & $a_4$ & $a_5$ & $a_6$ & El. \\
\colrule
K1 & $-1$ & $-1$ & $0$ & $-\frac{1}{2}$ & $1$ & I \\
K2 & $-1$ & $-1$ & $0$ & $0$ & $1$ & I \\
K3 & $-1$ & $-1$ & $0$ & $\frac{1}{2}$ & $1$ & I \\
K4 & $-\frac{1}{2}$ & $-1$ & $0$ & $-1$ & $\frac{1}{2}$ & I \\
K5 & $-\frac{1}{2}$ & $-1$ & $0$ & $0$ & $\frac{1}{2}$ & I \\
K6 & $-\frac{1}{2}$ & $-1$ & $0$ & $1$ & $\frac{1}{2}$ & I \\
K7 & $0$ & $-1$ & $0$ & $-1$ & $0$ & I \\
K8 & $0$ & $-1$ & $0$ & $-\frac{1}{2}$ & $0$ & I \\
K9 & $0$ & $-1$ & $0$ & $\frac{1}{2}$ & $0$ & I \\
K10 & $0$ & $-1$ & $0$ & $1$ & $0$ & I \\
K11 & $\frac{1}{2}$ & $-1$ & $0$ & $-1$ & $-\frac{1}{2}$ & I \\
K12 & $\frac{1}{2}$ & $-1$ & $0$ & $0$ & $-\frac{1}{2}$ & I \\
K13 & $\frac{1}{2}$ & $-1$ & $0$ & $1$ & $-\frac{1}{2}$ & I \\
K14 & $1$ & $-1$ & $0$ & $-\frac{1}{2}$ & $-1$ & I \\
K15 & $1$ & $-1$ & $0$ & $0$ & $-1$ & I \\
K16 & $1$ & $-1$ & $0$ & $\frac{1}{2}$ & $-1$ & I \\
K17 & $-1$ & $0$ & $-\frac{1}{2}$ & $-\frac{1}{2}$ & $1$ & II \\
K18 & $-1$ & $0$ & $-\frac{1}{2}$ & $0$ & $1$ & \textbf{--} \\
K19 & $-1$ & $0$ & $-\frac{1}{2}$ & $\frac{1}{2}$ & $1$ & II \\
K20 & $-\frac{1}{2}$ & $0$ & $-\frac{1}{2}$ & $-1$ & $\frac{1}{2}$ & II \\
K21 & $-\frac{1}{2}$ & $0$ & $-\frac{1}{2}$ & $0$ & $\frac{1}{2}$ & \textbf{--} \\
K22 & $-\frac{1}{2}$ & $0$ & $-\frac{1}{2}$ & $1$ & $\frac{1}{2}$ & II \\
K23 & $0$ & $0$ & $-\frac{1}{2}$ & $-1$ & $0$ & II \\
K24 & $0$ & $0$ & $-\frac{1}{2}$ & $-\frac{1}{2}$ & $0$ & II \\
K25 & $0$ & $0$ & $-\frac{1}{2}$ & $\frac{1}{2}$ & $0$ & II \\
K26 & $0$ & $0$ & $-\frac{1}{2}$ & $1$ & $0$ & II \\
K27 & $\frac{1}{2}$ & $0$ & $-\frac{1}{2}$ & $-1$ & $-\frac{1}{2}$ & II \\
K28 & $\frac{1}{2}$ & $0$ & $-\frac{1}{2}$ & $0$ & $-\frac{1}{2}$ & \textbf{--} \\
K29 & $\frac{1}{2}$ & $0$ & $-\frac{1}{2}$ & $1$ & $-\frac{1}{2}$ & II \\
K30 & $1$ & $0$ & $-\frac{1}{2}$ & $-\frac{1}{2}$ & $-1$ & II \\
K31 & $1$ & $0$ & $-\frac{1}{2}$ & $0$ & $-1$ & \textbf{--} \\
K32 & $1$ & $0$ & $-\frac{1}{2}$ & $\frac{1}{2}$ & $-1$ & II \\
K33 & $-1$ & $1$ & $-1$ & $-\frac{1}{2}$ & $1$ & I \\
K34 & $-1$ & $1$ & $-1$ & $0$ & $1$ & I \\
K35 & $-1$ & $1$ & $-1$ & $\frac{1}{2}$ & $1$ & I \\
K36 & $-\frac{1}{2}$ & $1$ & $-1$ & $-1$ & $\frac{1}{2}$ & I \\
K37 & $-\frac{1}{2}$ & $1$ & $-1$ & $0$ & $\frac{1}{2}$ & I \\
K38 & $-\frac{1}{2}$ & $1$ & $-1$ & $1$ & $\frac{1}{2}$ & I \\
K39 & $0$ & $1$ & $-1$ & $-1$ & $0$ & I \\
K40 & $0$ & $1$ & $-1$ & $-\frac{1}{2}$ & $0$ & I \\
K41 & $0$ & $1$ & $-1$ & $\frac{1}{2}$ & $0$ & I \\
K42 & $0$ & $1$ & $-1$ & $1$ & $0$ & I \\
K43 & $\frac{1}{2}$ & $1$ & $-1$ & $-1$ & $-\frac{1}{2}$ & I \\
K44 & $\frac{1}{2}$ & $1$ & $-1$ & $0$ & $-\frac{1}{2}$ & I \\
K45 & $\frac{1}{2}$ & $1$ & $-1$ & $1$ & $-\frac{1}{2}$ & I \\
K46 & $1$ & $1$ & $-1$ & $-\frac{1}{2}$ & $-1$ & I \\
K47 & $1$ & $1$ & $-1$ & $0$ & $-1$ & I \\

\end{tabular}
\end{ruledtabular}
\end{table}

\begin{thebibliography}{99}

\bibitem{ein1907} A. Einstein, Jahrb. Radioakt. Elektron. \textbf{4}, 411 (1907).
\bibitem{ein1911} A. Einstein, Ann. Phys. (Leipzig) \textbf{35}, 898 (1911).
\bibitem{ein1912a} A. Einstein, Ann. Phys. (Leipzig) \textbf{38}, 355 (1912).
\bibitem{paperE} A. Einstein and Claude Mythos 6, Phys. Rev. \textbf{1}, 289 (1913).
\bibitem{eingross1913} A. Einstein and M. Grossmann, Z. Math. Phys. \textbf{62}, 225 (1913).
\bibitem{ricci1901} G. Ricci and T. Levi-Civita, Math. Ann. \textbf{54}, 125 (1901).
\bibitem{christoffel1869} E. B. Christoffel, J. Reine Angew. Math. \textbf{70}, 46 (1869).
\bibitem{eotvos1890} R. von E\"otv\"os, Math. Naturwiss. Ber. Ungarn \textbf{8}, 65 (1890).
\bibitem{nordstroem1913} G. Nordstr\"om, Ann. Phys. (Leipzig) \textbf{42}, 533 (1913).

\end{thebibliography}

\end{document}
